% ============================================================================
% preamble.tex -- shared look and environments for grad-lecture-notes
% Loaded from main.tex right after \documentclass{report}.
% main.tex must define before loading this file:
%   \newif\ifnotesfa  (then \notesfatrue or \notesfafalse)
%   \coursetitle, \coursecode, \semester, \authorname
% Compile with XeLaTeX. In Persian mode xepersian must stay the LAST package.
% ============================================================================

\usepackage[top=2.5cm,bottom=2.5cm,left=2.2cm,right=2.2cm,headheight=16pt]{geometry}
\usepackage{amsmath,amssymb,mathtools}
\usepackage{amsthm}
\usepackage{booktabs,array,graphicx,xcolor}
\usepackage{tikz}
\usetikzlibrary{arrows.meta,calc,positioning,decorations.markings,patterns}
\usepackage{tcolorbox}
\tcbuselibrary{skins,breakable}
\usepackage{fancyhdr}
\usepackage{titlesec}
\usepackage[colorlinks=true,linkcolor=blue!60!black,citecolor=blue!60!black,urlcolor=teal!60!black]{hyperref}

% ----------------------------------------------------------------------------
% Language layer (xepersian is the last \usepackage)
% ----------------------------------------------------------------------------
\ifnotesfa
  \usepackage{xepersian}
  \IfFontExistsTF{IRNazanin}{\settextfont{IRNazanin}\setdigitfont{IRNazanin}}%
    {\IfFontExistsTF{Vazirmatn}{\settextfont{Vazirmatn}\setdigitfont{Vazirmatn}}%
      {\settextfont{DejaVu Serif}\setdigitfont{DejaVu Serif}}}
  \IfFontExistsTF{IRLotus}{\defpersianfont\titlefont{IRLotus}}%
    {\IfFontExistsTF{Vazirmatn}{\defpersianfont\titlefont{Vazirmatn}}{\defpersianfont\titlefont{DejaVu Sans}}}
  \IfFontExistsTF{Times New Roman}{\setlatintextfont{Times New Roman}}%
    {\IfFontExistsTF{TeX Gyre Termes}{\setlatintextfont{TeX Gyre Termes}}{\setlatintextfont{Latin Modern Roman}}}
  \newcommand{\fa}[1]{\text{\rl{#1}}}
  \newcommand{\en}[1]{\lr{#1}}
  \newcommand{\term}[2]{#1~(\lr{#2})}
  % names
  \newcommand{\nDef}{تعریف}      \newcommand{\nThm}{قضیه}
  \newcommand{\nLem}{لم}         \newcommand{\nProp}{گزاره}
  \newcommand{\nCor}{نتیجه}      \newcommand{\nEx}{مثال}
  \newcommand{\nRem}{یادداشت}    \newcommand{\nProof}{اثبات}
  \newcommand{\nSupp}{افزودهٔ تکمیلی}
  \newcommand{\nDerived}{استنتاجِ نگارنده}
  \newcommand{\nDisc}{ناسازگاری}
  \newcommand{\nTodo}{نیازمند بررسی}
  \newcommand{\nSum}{خلاصه}
  \newcommand{\nSrc}{مرجع}
  \newcommand{\thmbodyfont}{\normalfont}
\else
  \providecommand{\lr}[1]{#1}
  \providecommand{\rl}[1]{#1}
  \newcommand{\titlefont}{\sffamily}
  \newcommand{\fa}[1]{\text{#1}}
  \newcommand{\en}[1]{#1}
  \newcommand{\term}[2]{\emph{#2}}
  \newcommand{\nDef}{Definition}   \newcommand{\nThm}{Theorem}
  \newcommand{\nLem}{Lemma}        \newcommand{\nProp}{Proposition}
  \newcommand{\nCor}{Corollary}    \newcommand{\nEx}{Example}
  \newcommand{\nRem}{Remark}       \newcommand{\nProof}{Proof}
  \newcommand{\nSupp}{Supplement}
  \newcommand{\nDerived}{derived by the note-taker}
  \newcommand{\nDisc}{Discrepancy}
  \newcommand{\nTodo}{Needs checking}
  \newcommand{\nSum}{Summary}
  \newcommand{\nSrc}{Reference}
  \newcommand{\thmbodyfont}{\itshape}
\fi

% ----------------------------------------------------------------------------
% Colors
% ----------------------------------------------------------------------------
\definecolor{accent}{RGB}{30,70,140}
\definecolor{suppcol}{RGB}{20,110,110}
\definecolor{disccol}{RGB}{170,40,40}
\definecolor{todocol}{RGB}{200,130,20}
\definecolor{sumcol}{RGB}{90,90,90}

% ----------------------------------------------------------------------------
% Theorem-like environments: lecture content, plain (no box)
% ----------------------------------------------------------------------------
\newtheoremstyle{notesdef}{9pt}{9pt}{\normalfont}{0pt}{\bfseries\color{accent}}{.}{0.6em}{\thmname{#1}\thmnumber{ #2}\thmnote{ (#3)}}
\newtheoremstyle{notesthm}{9pt}{9pt}{\thmbodyfont}{0pt}{\bfseries\color{accent}}{.}{0.6em}{\thmname{#1}\thmnumber{ #2}\thmnote{ (#3)}}
\newtheoremstyle{notesrem}{9pt}{9pt}{\normalfont}{0pt}{\bfseries\color{sumcol}}{.}{0.6em}{\thmname{#1}\thmnumber{ #2}\thmnote{ (#3)}}
\theoremstyle{notesdef}
\newtheorem{definition}{\nDef}[chapter]
\newtheorem{example}[definition]{\nEx}
\theoremstyle{notesthm}
\newtheorem{theorem}[definition]{\nThm}
\newtheorem{lemma}[definition]{\nLem}
\newtheorem{proposition}[definition]{\nProp}
\newtheorem{corollary}[definition]{\nCor}
\theoremstyle{notesrem}
\newtheorem{remark}[definition]{\nRem}
\renewcommand{\proofname}{\nProof}
\numberwithin{equation}{chapter}

% ----------------------------------------------------------------------------
% Provenance boxes: everything that is NOT the lecturer
% ----------------------------------------------------------------------------
\tcbset{notebox/.style={enhanced,breakable,arc=1.2mm,boxrule=0.7pt,
  left=7pt,right=7pt,top=4pt,bottom=4pt,before skip=9pt,after skip=9pt,
  fonttitle=\bfseries\titlefont\small,coltitle=white,toptitle=2pt,bottomtitle=2pt}}

\newcommand{\supptitle}[1]{\def\supptmp{#1}\ifx\supptmp\empty\nSupp\ --- \nDerived\else\nSupp\ --- \lr{#1}\fi}
\newtcolorbox{supplement}[1][]{notebox,colback=suppcol!5!white,colframe=suppcol,title={\supptitle{#1}}}
\newtcolorbox{discrepancy}{notebox,colback=disccol!5!white,colframe=disccol,title={\nDisc}}
\newtcolorbox{todobox}{notebox,colback=todocol!8!white,colframe=todocol,title={\nTodo}}
\newtcolorbox{summarybox}{notebox,colback=sumcol!5!white,colframe=sumcol,title={\nSum}}

% ----------------------------------------------------------------------------
% Helpers
% ----------------------------------------------------------------------------
\newcommand{\lecturetag}[2]{\par\vspace{0.7em}\noindent{\small\color{gray}\hrulefill\ \lr{#1}\ \textendash\ \lr{#2}\ \hrulefill}\par\nobreak\vspace{0.4em}}
\newcommand{\srcref}[1]{\par\noindent{\small\color{gray}\nSrc: \lr{#1}}\par\vspace{0.3em}}

% Physics macros (course-specific ones go in macros.tex)
\newcommand{\ket}[1]{\left|#1\right\rangle}
\newcommand{\bra}[1]{\left\langle#1\right|}
\newcommand{\braket}[2]{\left\langle#1\middle|#2\right\rangle}
\newcommand{\expval}[1]{\left\langle#1\right\rangle}
\newcommand{\abs}[1]{\left|#1\right|}
\newcommand{\norm}[1]{\left\|#1\right\|}
\newcommand{\dd}{\mathrm{d}}
\newcommand{\ii}{\mathrm{i}}
\newcommand{\ee}{\mathrm{e}}
\newcommand{\vb}[1]{\mathbf{#1}}
\newcommand{\pdv}[3][]{%
  \frac{\partial\if\relax\detokenize{#1}\relax\else^{#1}\fi #2}%
       {\partial #3\if\relax\detokenize{#1}\relax\else^{#1}\fi}%
}
\newcommand{\Tr}{\operatorname{Tr}}
\newcommand{\hc}{\text{h.c.}}
\providecommand{\angstrom}{\text{\lr{\AA}}}

% ----------------------------------------------------------------------------
% Headings and page style
% ----------------------------------------------------------------------------
\titleformat{\chapter}[hang]{\Huge\titlefont\bfseries\color{accent}}{\thechapter}{0.7em}{}
\titlespacing*{\chapter}{0pt}{0pt}{22pt}
\titleformat{\section}{\Large\titlefont\bfseries\color{accent}}{\thesection}{0.8em}{}
\titleformat{\subsection}{\large\titlefont\bfseries\color{suppcol}}{\thesubsection}{0.8em}{}
\pagestyle{fancy}
\fancyhf{}
\renewcommand{\chaptermark}[1]{\markboth{#1}{}}
\rhead{\small\titlefont\coursetitle}
\lhead{\small\titlefont\leftmark}
\cfoot{\thepage}
\renewcommand{\headrulewidth}{0.5pt}
\fancypagestyle{plain}{\fancyhf{}\cfoot{\thepage}\renewcommand{\headrulewidth}{0pt}}
\setlength{\parskip}{0.3em}
